\chapter*{Abstract}

%CONTEXT
%hybrid systems
%use floating point arithmetic, which causes errors
%thats bad for some applications
%STATEMENT
%lets use exact real arithmetic
%STATE OF THE ART
%drawing from previous research on computable reals, backed by domain theory, as well as a monadic formalization of hybridness
%GOALS
%implement a hybrid simulator and assess its performance
%METHODS
%in Haskell
%using ERA libraries
%benchmarking
%FINDINGS
%ERA overall is much slower than floating'point, as expected
%some strategies are better than others
%IMPLICATIONS
%the practical tool, Jaguar (namedrop sooner) is pretty cool and can be used for stuff
%future ERA libraries can also learn from these results

Like most programs, simulations of hybrid systems mostly use floating-point arithmetic, which can introduce errors and lead to false conclusions, especially in sensitive and critical systems. This dissertation aims to fix this problem through the use of exact real arithmetic in hybrid programs. Drawing from previous research on computable real numbers backed by domain theory, as well as a monadic formalization of hybridness, a new hybrid simulator Jaguar is developed in Haskell using preexisting exact real arithmetic libraries. The greatest challenges in development are the semidecidability of comparison between computable reals and the implementation of a differential equation solver. As expected, the final product is overall slower than equivalent programs using floating-point; but there are interesting performance differences between various exact real arithmetic libraries due to their different strategies and characteristics, as explored in the benchmarking. These results are doubtlessly of interest from the perspective of designing new performant exact real arithmetic libraries, and Jaguar as a standalone tool has applications in (small) physical simulations, formal verification and more.

\paragraph{Keywords} exact real arithmetic, hybrid systems, Jaguar, computable real numbers

\cleardoublepage

\chapter*{Resumo}

Tal como a maioria dos programas, as simulações de sistemas híbridos utilizam maioritariamente aritmética de vírgula flutuante, o que pode introduzir erros e levar a conclusões incorretas, especialmente em sistemas críticos ou sensíveis. Esta dissertação procura resolver este problema através da utilização de aritmética real exata em programas híbridos. Com base em investigação existente sobre números reais computáveis baseada em teoria de domínios, bem como numa formalização monádica da natureza híbrida destes sistemas, foi desenvolvido um novo simulador híbrido em Haskell chamado Jaguar, recorrendo a bibliotecas pré-existentes de aritmética real exata. Os maiores desafios durante o desenvolvimento foram a semidecidibilidade da comparação entre números reais computáveis e a implementação de um solver de equações diferenciais. Como seria de esperar, o produto final é, de um modo geral, mais lento do que programas equivalentes que usam aritmética de vírgula flutuante. No entanto, as bibliotecas de aritmética real exata apresentam entre si diferenças de desempenho, resultantes das suas diferentes estratégias e características, tal como demonstrado no benchmarking. Estes resultados são relevantes para o desenvolvimento de futuras bibliotecas de aritmética real exata de alto desempenho. Além disso, o Jaguar, enquanto ferramenta autónoma, tem aplicações em simulações físicas de pequena escala, verificação formal e outras áreas.

\paragraph{Palavras-chave} aritmética real exata, sistemas híbridos, Jaguar, números reais computáveis

\cleardoublepage
